\begin{abstract}
The success of Transformers lies in their ability to improve inference through two complementary strategies: the permanent refinement of model parameters via \textit{in-weight learning} (IWL), and the ephemeral modulation of inferences via \textit{in-context learning} (ICL), which leverages contextual information maintained in the model's activations.
Evolutionary biology tells us that the predictability of the environment across timescales predicts the extent to which analogous strategies should be preferred. Genetic \textit{evolution} adapts to stable environmental features by gradually modifying the genotype over generations. Conversely, environmental volatility favors \textit{plasticity}, which enables a single genotype to express different traits within a lifetime, provided there are reliable cues to guide the adaptation.
We operationalize these dimensions (environmental stability and cue reliability) in controlled task settings (sinusoid regression and Omniglot classification) to characterize their influence on learning in Transformers.
We find that stable environments favor IWL, often exhibiting a sharp transition when conditions are static. Conversely, reliable cues favor ICL, particularly when the environment is volatile.
Furthermore, an analysis of learning dynamics reveals task-dependent transitions between strategies (ICL $\to$ IWL and vice versa). We demonstrate that these transitions are governed by (1) the asymptotic optimality of the strategy with respect to the environment, and (2) the optimization cost of acquiring that strategy, which depends on the task structure and the learner's inductive bias.
\end{abstract}

% \begin{abstract}
% Transformer models learn in two distinct modes: in-weights learning (IWL), encoding knowledge into model weights, and in-context learning (ICL), adapting flexibly to context without weight modification.
% To better understand the interplay between these learning modes, we draw inspiration from evolutionary biology's analogous adaptive strategies: genetic encoding (akin to IWL, adapting over generations and fixed within an individual's lifetime) and phenotypic plasticity (akin to ICL, enabling flexible behavioral responses to environmental cues). In evolutionary biology, environmental predictability dictates the balance between these strategies: stability favors genetic encoding, while reliable predictive cues promote phenotypic plasticity.
% We experimentally operationalize these dimensions of predictability and systematically investigate their influence on the ICL/IWL balance in Transformers. Using regression and classification tasks, we show that high environmental stability decisively favors IWL, as predicted, with a sharp transition at maximal stability. Conversely, high cue reliability enhances ICL efficacy, particularly when stability is low. Furthermore, learning dynamics reveal task-contingent temporal evolution: while a canonical ICL-to-IWL shift occurs in some settings (e.g., classification with many classes), we demonstrate that scenarios with easier IWL (e.g., fewer classes) or slower ICL acquisition (e.g., regression) can exhibit an initial IWL phase later yielding to ICL dominance. These findings support a relative-cost hypothesis for explaining these learning mode transitions, establishing predictability as a critical factor governing adaptive strategies in Transformers, and offering novel insights for understanding ICL and guiding training methodologies.
% \end{abstract}

% \begin{abstract}
% Transformer models can learn in two quite different ways: in-weights learning (IWL; encoding knowledge in model weights) vs. in-context learning (ICL; flexibly adapting to context without weight changes). To better understand the balance and dynamics between these two types of learning, we turn to evolutionary biology, which exhibits analogous modes of adaptation: genetic encoding of behavior (akin to IWL; adapting over longer timescales but fixed within the lifetime of an individual) vs. phenotypic plasticity (akin to ICL; where behavior is not hardcoded and instead flexibly adapts to the environment). In evolutionary biology, the balance between these two modes of adaptation depends on environmental predictability; specifically, environmental stability favors genetic encoding, while the reliability of cues in predicting future conditions favors phenotypic plasticity. In this work, we experimentally operationalize these two types of predictability, and systematically investigate how they govern the ICL/IWL balance in Transformers. Using simple regression and classification tasks, we show that high environmental stability strongly promotes IWL over ICL, with a sharp transition to IWL at maximal stability, while high cue reliability enhances ICL efficacy, especially when environmental stability is low. Furthermore, learning dynamics reveal that the time course of ICL and IWL depends on the task: the typical shift from ICL to IWL occurs in some settings (e.g., classification with many classes), but we demonstrate that in settings where IWL is easier (e.g., classification with fewer classes) or when ICL is more slowly acquired (e.g., regression) learning can show an initial IWL phase later ceding to ICL dominance. These findings support a relative-cost hypothesis for shifts between types of learning and establish predictability as a critical factor shaping adaptive learning in Transformers, offering insights for understanding ICL and guiding training methodologies.
% \end{abstract}